\section*{Bab \ref{ch:b1:multivar} --- Fungsi Dua Variabel}

\begin{solution}{exo:b1:multivar:1}
$\dfrac{\partial f}{\partial x} = 2xy + y\,\eu^{xy}$,
$\dfrac{\partial f}{\partial y} = x^2 + x\,\eu^{xy}$.

$\dfrac{\partial g}{\partial x} = \dfrac{2x}{x^2+y^2}$,
$\dfrac{\partial g}{\partial y} = \dfrac{2y}{x^2+y^2}$.

$\dfrac{\partial h}{\partial x} = \dfrac{-y/x^2}{1 + y^2/x^2} =
\dfrac{-y}{x^2+y^2}$,
$\dfrac{\partial h}{\partial y} = \dfrac{x}{x^2+y^2}$.
\end{solution}

\begin{solution}{exo:b1:multivar:2}
Dalam \omterm{prop:b1:vspaces:coordinates}{koordinat} kutub ($\rho \to 0$):

$\abs{f} = \dfrac{\rho^3\abs{\cos^2\theta\sin\theta}}{\rho^2} \leq
\rho \to 0$: jadi \omterm{def:b1:continuity:continuous}{kontinu}.

$g = \cos\theta\sin\theta$: yang tak bergantung pada $\rho$, dan mengambil nilai
yang berbeda sepanjang sinar garis yang berbeda (bandingkanlah
\cref{ex:b1:multivar:trap}): jadi tak berlimit, dan tak \omterm{def:b1:continuity:continuous}{kontinu}.

$\abs h \leq \dfrac{\rho^3(\abs{\cos^3\theta} +
\abs{\sin^3\theta})}{\rho^2} \leq 2\rho \to 0$: jadi \omterm{def:b1:continuity:continuous}{kontinu}.
\end{solution}

\begin{solution}{exo:b1:multivar:3}
Sepanjang $y = mx$: $f(x, mx) = \dfrac{m x^3}{x^4 + m^2 x^2} =
\dfrac{mx}{x^2 + m^2} \to 0$ (untuk $m \neq 0$; sedangkan sepanjang $y = 0$ dan
sumbu $y$, $f = 0$). Jadi setiap limit garis lurusnya $0$. Tetapi pada
parabola $y = x^2$:
\[
f(x, x^2) = \frac{x^4}{x^4 + x^4} = \frac12 :
\]
jadi barisan $\bigl(\frac1n, \frac{1}{n^2}\bigr) \to (0,0)$ mempunyai
$f \to \frac12 \neq 0$. Jadi tak \omterm{def:b1:continuity:continuous}{kontinu}: karena garis tak mencukupi untuk
menguji limit dua variabel.
\end{solution}

\begin{solution}{exo:b1:multivar:4}
$\frac{\partial f}{\partial x} = 3x^2y^2 + y\cos(xy)$; lalu
\[
\frac{\partial^2 f}{\partial y\,\partial x}
= 6x^2 y + \cos(xy) - xy\sin(xy) .
\]
$\frac{\partial f}{\partial y} = 2x^3 y + x\cos(xy)$; lalu
\[
\frac{\partial^2 f}{\partial x\,\partial y}
= 6x^2 y + \cos(xy) - xy\sin(xy) :
\]
jadi sama, sebagaimana dijanjikan Schwarz.
\end{solution}

\begin{solution}{exo:b1:multivar:5}
Menurut \cref{thm:b1:multivar:chain} dengan $(x, y) = (\cos t, \sin t)$:
\[
g'(t) = -\sin t\,\frac{\partial f}{\partial x}(\cos t, \sin t)
+ \cos t\,\frac{\partial f}{\partial y}(\cos t, \sin t)
= \bigl\langle \nabla f,\ (-\sin t, \cos t)\bigr\rangle .
\]
Pada sebuah ekstremum $g$, $g'(t) = 0$: sehingga $\nabla f$ \omterm{def:b1:euclid:orthogonal}{ortogonal} terhadap
vektor singgung $(-\sin t, \cos t)$ lingkarannya, jadi sejajar
vektor jari-jarinya $(\cos t, \sin t)$ (karena \omterm{def:b1:euclid:orthogonal}{pelengkap ortogonal}
bidang atas sebuah vektor satuan adalah garis yang direntangnya, yang diambil
tegak lurus). Inilah contoh tersederhana sebuah pengganda
Lagrange.
\end{solution}

\begin{solution}{exo:b1:multivar:6}
$f$: $\nabla f = (2x + y - 3,\; x + 2y) = 0$: jadi $y = -\frac x2$ dan
$2x - \frac x2 = 3$: sehingga $x = 2$, $y = -1$. Orde duanya: $r = 2$, $s =
1$, $t = 2$: $rt - s^2 = 3 > 0$, $r > 0$: jadi minimum lokal (bahkan global ---
karena kuadratik) di $(2, -1)$, bernilai $f(2,-1) = -3$.

$g$: $\nabla g = (2x,\; -2y + 4) = 0$: jadi titik $(0, 2)$; $r = 2$, $s
= 0$, $t = -2$: $rt - s^2 = -4 < 0$: jadi pelana.

$h$: $\nabla h = \bigl(4x^3 - 4(x - y),\; 4y^3 + 4(x-y)\bigr) = 0$.
Menambahkannya: $x^3 + y^3 = 0$, sehingga $y = -x$; lalu menyulihkannya: $4x^3 - 8x = 0$:
jadi $x \in \{0, \pm\sqrt2\}$. Titik kritisnya: $(0,0)$,
$(\sqrt2, -\sqrt2)$, $(-\sqrt2, \sqrt2)$.
Di $(\pm\sqrt2, \mp\sqrt2)$: $r = 12x^2 - 4 = 20$, $s = 4$, $t =
20$: $rt - s^2 > 0$, $r > 0$: jadi minimum lokal (bernilai $h = 4 + 4 - 16 =
-8$). Di $(0,0)$: $r = t = -4$, $s = 4$: $rt - s^2 = 0$: jadi bungkam.
Periksalah: $h(x, x) = 2x^4 > 0$ dan $h(x, -x) = 2x^4 - 8x^2 < 0$ bagi
$x \neq 0$ yang kecil: jadi kedua tandanya ada di setiap \omterm{def:b1:topology:open}{persekitarannya} --- sehingga tak ada
ekstremum di titik asalnya.
\end{solution}

\begin{solution}{exo:b1:multivar:7}
$x^2 - y^2$: \omterm{def:b1:multivar:partial}{turunan parsial} keduanya $2$ dan $-2$: jadi berjumlah $0$. $xy$: kedua
\omterm{def:b1:multivar:partial}{turunan parsial} kedua yang murninya lenyap. $\eu^x\cos y$: $\frac{\partial^2}{\partial
x^2} = \eu^x\cos y$, $\frac{\partial^2}{\partial y^2} =
-\eu^x\cos y$: jadi berjumlah $0$. $\ln(x^2+y^2)$: dari
\cref{exo:b1:multivar:1},
\[
\frac{\partial^2}{\partial x^2}\ln(x^2+y^2)
= \frac{2(x^2+y^2) - 4x^2}{(x^2+y^2)^2}
= \frac{2(y^2 - x^2)}{(x^2+y^2)^2},
\]
dan versi $y$-nya adalah lawannya: jadi berjumlah $0$.
\end{solution}

\begin{solution}{exo:b1:multivar:8}
\emph{\omterm{def:b1:linmaps:projection}{Proyeksinya}:} bidang $P$ bernormal $n = (1, 2, -1)$ dan
lewat $A = (4, 0, 0)$. Adapun titik yang terdekat dengan titik asalnya adalah
$O$ yang diproyeksikan: $p = O + \frac{\langle A - O, n\rangle}{\norm
n^2}\,n = \frac{4}{6}(1,2,-1) = \bigl(\frac23, \frac43,
-\frac23\bigr)$, pada jarak $\frac{4}{\sqrt 6}$. (Rumusnya: jarak
dari titik asalnya ke bidang $\langle X, n \rangle = c$ adalah
$\frac{\abs c}{\norm n}$ dengan $c = 4$.)

\emph{Peminimalannya:} $f(x,y) = x^2 + y^2 + (x + 2y - 4)^2$. Dengan menulis
$w = x + 2y - 4$ bagi faktor terakhirnya:
\[
\nabla f = \bigl(2x + 2w,\; 2y + 4w\bigr) = 0
\iff x = -w \text{ dan } y = -2w .
\]
Menyulihkannya ke dalam definisi $w$: $w = -w - 4w - 4$, sehingga $w =
-\frac23$, yang memberikan $x = \frac23$, $y = \frac43$, dan $z = w =
-\frac23$. Jadi titik yang sama seperti lewat \omterm{def:b1:linmaps:projection}{proyeksinya}, pada jarak
$\sqrt{\frac49 + \frac{16}{9} + \frac49} = \frac{2\sqrt 6}{3} =
\frac{4}{\sqrt6}$; dan ia minimum karena $f \to \infty$ di takhingga
(yaitu kuadratik yang tegas positif ditambah suku \omterm{def:b1:linmaps:def}{linear}).
\end{solution}

\begin{solution}{exo:b1:multivar:9}
Dengan $\rho = \sqrt{x^2+y^2}$: $\frac{\partial \rho}{\partial x} =
\frac{x}{\rho}$, sehingga
\[
\frac{\partial f}{\partial x} = \varphi'(\rho)\,\frac{x}{\rho},
\qquad
\frac{\partial^2 f}{\partial x^2}
= \varphi''(\rho)\,\frac{x^2}{\rho^2}
+ \varphi'(\rho)\,\Bigl(\frac{1}{\rho} -
\frac{x^2}{\rho^3}\Bigr),
\]
(lewat kaidah hasil bagi dan kaidah rantai). Menambahkan ungkapan $y$ yang setangkup:
maka suku $\varphi''$-nya menghimpun $\frac{x^2 + y^2}{\rho^2} = 1$, sedangkan
suku $\varphi'$-nya menghimpun $\frac{2}{\rho} - \frac{x^2 +
y^2}{\rho^3} = \frac{1}{\rho}$:
\[
\Delta f = \varphi'' + \frac{\varphi'}{\rho} .
\]
Harmonik yang radial: pecahkanlah $\varphi'' + \frac{\varphi'}{\rho} = 0$: karena
persamaan orde pertamanya $u' + \frac u\rho = 0$ bagi $u = \varphi'$
memberikan $u = \frac{c}{\rho}$ (\cref{thm:b1:diffeq:homogeneous1}),
lalu $\varphi = c\ln\rho + d$. Jadi fungsi harmonik yang radial adalah
$c\,\ln\sqrt{x^2 + y^2} + d$ --- yaitu potensial logaritmik
\cref{exo:b1:multivar:7} dan konstantanya, tak ada yang lain.
\end{solution}

\begin{solution}{exo:b1:multivar:10}
Untuk $z = xy$ di $(1,1)$: \omterm{def:b1:multivar:partial}{turunan parsialnya} $y = 1$ dan $x = 1$, sehingga bidang
singgungnya $z = 1 + (x - 1) + (y - 1) = x + y - 1$. Adapun bidang singgung
yang mendatar berarti kedua \omterm{def:b1:multivar:partial}{turunan parsialnya} lenyap, yaitu $\nabla f = 0$:
jadi menurut \cref{ex:b1:multivar:extrema}, tepat titik kritisnya
$(0, 0)$ dan $(1, 1)$, dengan bidang mendatarnya $z = 0$ dan $z =
-1$. Sehingga ``bidang singgung mendatar'' dan ``\omterm{thm:b1:multivar:critical}{titik kritis}'' merupakan
gagasan yang sama, yang dilihat pada grafiknya dan pada rumusnya.
\end{solution}

\begin{solution}{exo:b1:multivar:11}
\begin{enumerate}
  \item Bagi $y$ yang tetap, fungsi satu variabel $x \mapsto
        f(x,y)$ berturunan nol pada $\R$, jadi tetap
        (\cref{thm:b1:derivative:mvt}): sehingga $f(x, y) = f(0, y)$ untuk
        semua $x$: jadi $f$ hanya bergantung pada $y$.
  \item Misalkan $g(u, v) = f\bigl(\frac{u+v}2, \frac{u-v}2\bigr)$.
        Menurut kaidah rantainya (\cref{thm:b1:multivar:chain}, yang diterapkan
        pada variabel $v$ dengan $u$ yang dibekukan):
        \[
        \frac{\partial g}{\partial v}
        = \frac12\,\frac{\partial f}{\partial x}
        - \frac12\,\frac{\partial f}{\partial y} = 0 .
        \]
        Menurut (1), $g$ hanya bergantung pada $u$: $g(u, v) = \varphi(u)$
        dengan $\varphi$ yang berkelas $C^1$, dan dengan membalikkan
        perubahan variabelnya ($u = x + y$, $v = x - y$),
        \[
        f(x, y) = \varphi(x + y).
        \]
        Sebaliknya setiap $f$ yang demikian memenuhi $f_x = f_y =
        \varphi'$: jadi penyelesaiannya tepat fungsi $C^1$
        atas $x + y$.
\end{enumerate}
\end{solution}

\begin{solution}{exo:b1:multivar:12}
Pada cakram \omterm{def:b1:topology:open}{terbukanya}, sebuah ekstremum tentu kritis: $\nabla f = (y,
x) = 0$ hanya di $(0,0)$, tempat $f = 0$; dan ia sebuah pelana ($f(\pm
\varepsilon, \pm\varepsilon) = \varepsilon^2 > 0 > -\varepsilon^2
= f(\pm\varepsilon, \mp\varepsilon)$), sehingga tak ada ekstremum di situ. Menurut
teorema nilai ekstrem yang diakui itu batasnya tercapai,
yang tentu pada lingkaran batasnya. Di situ, bersama
\cref{exo:b1:multivar:5},
\[
f(\cos t, \sin t) = \cos t\sin t = \frac{\sin 2t}{2}
\in \intcc{-\tfrac12}{\tfrac12},
\]
dengan maksimum $\frac12$ di $t = \frac\pi4, \frac{5\pi}4$ (yaitu titik
$\pm\frac{1}{\sqrt2}(1,1)$) dan minimum $-\frac12$ di $t =
\frac{3\pi}4, \frac{7\pi}4$ (yaitu titik $\pm\frac{1}{\sqrt2}(1,-1)$).
Jadi maksimum globalnya $\frac12$, dan minimum globalnya $-\frac12$.
\end{solution}

\begin{solution}{pb:b1:multivar:1}
\textbf{1.} Menjabarkan ruas kanannya:
\[
r\Bigl(h + \frac srk\Bigr)^{\!2} + \frac{rt - s^2}{r}k^2
= rh^2 + 2shk + \frac{s^2}{r}k^2 + \frac{rt - s^2}{r}k^2
= q(h,k) .
\]
Jika $rt - s^2 > 0$: maka kedua kuadratnya membawa tanda faktor $r$,
dan $q(h,k) = 0$ memaksa $k = 0$ lalu $h = 0$: jadi bertanda tegas seperti $r$
di luar titik asalnya. Jika $rt - s^2 < 0$: maka $q(1, 0) = r$ sedangkan
$q\bigl(-\frac sr, 1\bigr) = \frac{rt - s^2}{r}$ bertanda
berlawanan: jadi kedua tandanya terjadi.

\textbf{2.} Jika $r = 0 \neq t$, tukarlah peran $h$ dan $k$
(yaitu kuadrat lengkapnya pada $k$), dengan kesimpulan yang sama; dan perhatikanlah
$rt - s^2 = -s^2 < 0$ segera setelah $s \neq 0$, dan memang $q =
2shk + tk^2$ lalu mengambil kedua tandanya ($k$ yang kecil terhadap $h$). Jika
$r = t = 0$: $q = 2shk$, yang berkedua tanda jika dan hanya jika $s \neq 0$, yaitu jika dan hanya jika
$rt - s^2 = -s^2 < 0$. Ringkasnya: $q$ mengambil kedua tandanya $\iff rt -
s^2 < 0$; sedangkan $q$ lenyap hanya di titik asalnya (jadi tegas) $\iff rt -
s^2 > 0$ (yang memaksa $r \neq 0$, karena $r = 0$ memberikan $q(1,0) =
0$).

\textbf{3.} Dengan $f_x = rx + sy + \beta$ dan $f_y = sx + ty +
\gamma$, penjabaran langsungnya memberikan
\[
f(a + h, b + k) = f(a, b) + h\,f_x(a,b) + k\,f_y(a,b)
+ \tfrac12 q(h, k),
\]
tanpa suku yang lebih tinggi (karena fungsinya \omterm{def:b1:poly:def}{polinomial} berderajat
$2$). Pada sebuah \omterm{thm:b1:multivar:critical}{titik kritis} bagian \omterm{def:b1:linmaps:def}{linearnya} lenyap: $f(X_0 + H)
- f(X_0) = \frac12 q(H)$ \emph{dengan persis}, sehingga kajian tanda
pertanyaan 1--2 menggolongkannya: yaitu minimum global yang tegas ($rt - s^2 >
0$, $r > 0$), maksimum global yang tegas ($rt - s^2 > 0$, $r < 0$),
dan pelana --- yakni kedua tandanya di setiap \omterm{def:b1:topology:open}{persekitarannya} --- ketika $rt - s^2 <
0$. Ini membuktikan \cref{met:b1:multivar:monge} bagi fungsi
kuadratik.

\textbf{4.} Bentuk $q - c(h^2 + k^2)$ berdata $r - c$, $s$,
$t - c$. Dengan $c = \frac{rt - s^2}{r + t}$ (perhatikanlah $t > 0$ karena
$rt > s^2 \geq 0$ dan $r > 0$):
\[
(r - c)(t - c) - s^2 = rt - s^2 - c(r + t) + c^2 = c^2 \geq 0,
\]
dan $r - c \geq 0$ (karena $c \leq r \iff rt - s^2 \leq r^2 + rt$,
yang benar). Jika $r - c > 0$, maka kuadrat lengkap pertanyaan 1 menunjukkan
$q - c(h^2 + k^2) \geq 0$; sedangkan jika $r - c = 0$, maka $s = 0$ (karena dari
$-s^2 = -c^2 \leq \dots$ yang memaksa besaran yang ditampilkan
$\geq 0$ dengan faktor pertamanya $0$) dan bentuknya menjadi $(t - c)k^2 \geq
0$. Jadi pada kedua kasusnya $q(h,k) \geq c(h^2 + k^2)$.

\textbf{5.} Menurut pertanyaan 3--4, $f(X) = f(X_0) + \frac12 q(X -
X_0) \geq f(X_0) + \frac c2\norm{X - X_0}^2 \to +\infty$ ketika
$\norm X \to \infty$: jadi $f$ bersifat koersif, dan ketaksamaannya
tegas untuk $X \neq X_0$: sehingga titik kritisnya pemberi minimum global
yang tunggal. Adapun bagi kubik $f = x^3 + y^3 - 3xy$ pada
\cref{ex:b1:multivar:extrema}, tak ada kesimpulan demikian yang berlaku: karena $f(x,0)
= x^3 \to -\infty$, dan minimum lokalnya di $(1,1)$ bukan
yang global --- sebab kepersisan penjabaran kuadratiknyalah yang gagal.

\textbf{6.} Menjabarkan kuadrat normanya:
\[
f(u,v) = u^2\norm{C_1}^2 + 2uv\,\langle C_1, C_2\rangle +
v^2\norm{C_2}^2 - 2u\,\langle C_1, b\rangle - 2v\,\langle C_2,
b\rangle + \norm b^2 ,
\]
yaitu fungsi kuadratik atas $(u, v)$ yang bagian kuadratiknya
$\frac12(ru^2 + 2suv + tv^2)$ dengan $r = 2\norm{C_1}^2$, $s =
2\langle C_1, C_2\rangle$, $t = 2\norm{C_2}^2$.

\textbf{7.} $rt - s^2 = 4\bigl(\norm{C_1}^2\norm{C_2}^2 -
\langle C_1, C_2\rangle^2\bigr) \geq 0$ menurut Cauchy--Schwarz
(\cref{thm:b1:euclid:cs}), dengan kesamaannya tepat ketika $C_1, C_2$
sebanding (atau salah satunya nol), yaitu ketika pasangannya
terkait. Jadi $rt - s^2 > 0 \iff (C_1, C_2)$ \omterm{def:b1:vspaces:free}{bebas}.

\textbf{8.} $\nabla f = 0$ terbaca
\[
\norm{C_1}^2 u + \langle C_1, C_2\rangle v = \langle C_1,
b\rangle,
\qquad
\langle C_1, C_2\rangle u + \norm{C_2}^2 v = \langle C_2,
b\rangle,
\]
yaitu \omterm{pb:b1:multivar:1}{persamaan normalnya} dengan matriks Gram di kirinya; dan ia berkata
$\langle uC_1 + vC_2 - b,\ C_i\rangle = 0$ untuk $i = 1, 2$, yaitu
$b - p \perp \operatorname{Vect}(C_1, C_2)$ bagi $p = uC_1 +
vC_2$. Menurut Bagian I (pertanyaan 3, 5) \omterm{thm:b1:multivar:critical}{titik kritis} yang tunggal itu
pemberi minimum global yang tunggal, dan menurut
\cref{thm:b1:euclid:projection} penciriannya ``$p \in
F$, $b - p \perp F$'' mengenali $p$ sebagai \omterm{thm:b1:euclid:projection}{proyeksi
ortogonal} $b$ ke $F = \operatorname{Vect}(C_1, C_2)$.

\textbf{9.} $\norm{b - p}^2 = \norm b^2 - 2\langle b, p\rangle +
\norm p^2$, dan $\langle p, b - p\rangle = 0$ memberikan $\norm p^2 =
\langle p, b\rangle$: sehingga minimumnya sama dengan $\norm b^2 - \langle p,
b\rangle$. Pythagoras: $\norm b^2 = \norm p^2 + \norm{b - p}^2$
--- jadi vektor datanya terbelah menjadi bagian yang terjelaskan $p$ dan
bagian sisaannya $b - p$, yang saling \omterm{def:b1:euclid:orthogonal}{ortogonal}.

\textbf{10.} $E(\alpha, \beta) = \norm{\alpha C_1 + \beta C_2 -
b}^2$ dengan $C_1, C_2, b$ yang dinyatakan itu; sehingga \omterm{pb:b1:multivar:1}{persamaan normal}
pertanyaan 8 adalah, entri demi entri,
\[
n\alpha + \Bigl(\sum x_i\Bigr)\beta = \sum y_i,
\qquad
\Bigl(\sum x_i\Bigr)\alpha + \Bigl(\sum x_i^2\Bigr)\beta = \sum
x_i y_i .
\]

\textbf{11.} \omterm{def:b1:arith:divides}{Membaginya} dengan $n$: $\alpha + \beta\bar x = \bar y$
dan $\alpha\bar x + \beta\bigl(v_x + \bar x^2\bigr) = c_{xy} +
\bar x\bar y$. Lalu menyulihkan $\alpha = \bar y - \beta\bar x$ ke dalam
yang kedua: $\beta v_x = c_{xy}$, sehingga
\[
\beta = \frac{c_{xy}}{v_x}, \qquad \alpha = \bar y - \beta\bar
x ,
\]
dan persamaan pertamanya berkata persis bahwa $(\bar x, \bar y)$
terletak pada garisnya.

\textbf{12.} $v_x = \frac1n\sum(x_i - \bar x)^2 \geq 0$, yang nol
jika dan hanya jika setiap $x_i$ sama dengan $\bar x$. Dan $rt - s^2 = 4\bigl(n\sum
x_i^2 - (\sum x_i)^2\bigr) = 4n^2 v_x$: jadi syarat kebebasan
pertanyaan 7 adalah $v_x > 0$, yaitu $x_i$ yang tak semuanya sama.

\textbf{13.} Dengan $\alpha = \bar y - \beta\bar x$, pusatkanlah
datanya ($\tilde x_i = x_i - \bar x$, $\tilde y_i = y_i - \bar y$):
\[
E(\alpha, \beta) = \sum_i(\tilde y_i - \beta\tilde x_i)^2
= n\,v_y - 2\beta\,n\,c_{xy} + \beta^2 n\,v_x ,
\]
yang minimal di $\beta = c_{xy}/v_x$ dengan nilai $E_{\min} =
n\bigl(v_y - \frac{c_{xy}^2}{v_x}\bigr) = n v_y(1 - \rho^2)$.
Karena $E_{\min} \geq 0$ dan $v_y > 0$: maka $\rho^2 \leq 1$, dan
$\abs\rho = 1$ jika dan hanya jika $E_{\min} = 0$, yaitu jika dan hanya jika setiap sisaannya
lenyap: sehingga titiknya terletak persis pada garisnya.

\textbf{14.} $n = 4$: $\bar x = \frac64 = 1.5$, $\bar y = 2$,
$\sum x_i^2 = 14$ sehingga $v_x = 3.5 - 2.25 = 1.25$, $\sum x_iy_i =
0 + 1 + 6 + 12 = 19$ sehingga $c_{xy} = 4.75 - 3 = 1.75$. Maka
\[
\beta = \frac{1.75}{1.25} = 1.4,
\qquad
\alpha = 2 - 1.4\times1.5 = -0.1 :
\qquad y = 1.4\,x - 0.1 .
\]
Nilai cocoknya $-0.1,\ 1.3,\ 2.7,\ 4.1$; sisaannya $0.1,\ -0.3,\
0.3,\ -0.1$; dan $E_{\min} = 0.01 + 0.09 + 0.09 + 0.01 = 0.2$
(periksa: $v_y = \frac{26}4 - 4 = 2.5$ dan $4(2.5 -
\frac{1.75^2}{1.25}) = 4(2.5 - 2.45) = 0.2$).

\textbf{15.} Kedua \omterm{pb:b1:multivar:1}{persamaan normal} pertanyaan 10 tepat
$\sum_i\varepsilon_i = 0$ dan $\sum_i x_i\varepsilon_i =
0$: jadi vektor sisaannya \omterm{def:b1:euclid:orthogonal}{ortogonal} terhadap $C_1 = (1, \dots, 1)$
dan terhadap $C_2 = (x_i)$ --- yakni terhadap seluruh ruang modelnya. Khususnya
garis yang optimal selalu menyeimbangkan galatnya: karena galatnya
berjumlah nol.

\textbf{16.} $\frac{\dd}{\dd\alpha}\sum(y_i - \alpha)^2 =
-2\sum(y_i - \alpha) = 0$ memberikan $\alpha = \bar y$ (dan
\omterm{def:b1:derivative:def}{turunan} keduanya $2n > 0$ menjadikannya minimum global, karena
fungsinya kuadratik yang koersif pada satu variabel). Adapun nilai
minimalnya $\sum(y_i - \bar y)^2 = n v_y$: jadi variansnya adalah
galat kuadratik yang tak tereduksikan pada model yang tetap.

\textbf{17.} $\frac{\dd}{\dd\beta}\sum(y_i - \beta x_i)^2 =
-2\sum x_i(y_i - \beta x_i) = 0$ memberikan $\beta_0 = \frac{\sum
x_iy_i}{\sum x_i^2}$. Dengan menulis kedua kemiringannya atas besaran
yang terpusat: $\beta_0 = \frac{c_{xy} + \bar x\bar y}{v_x + \bar
x^2}$ sama dengan $\frac{c_{xy}}{v_x}$ jika dan hanya jika $c_{xy}\bar x^2 = \bar
x\bar y\,v_x$, yaitu jika dan hanya jika $\bar x = 0$ (yakni absis yang terpusat) atau
$\bar y = \beta\bar x$ --- dengan yang terakhirnya berarti $\alpha = 0$: jadi
\omterm{pb:b1:multivar:1}{garis regresi} penuhnya sudah lewat titik asalnya.

\textbf{18.} $g(\tau) = \norm{M - P_0}^2 - 2\tau\langle M - P_0,
w\rangle + \tau^2$ minimal di $\tau^* = \langle M - P_0,
w\rangle$, dengan nilai $d(M,D)^2 = \norm{M - P_0}^2 - \langle M -
P_0, w\rangle^2$. Pada bidangnya, lengkapkanlah $w$ menjadi \omterm{def:b1:vspaces:free}{basis}
\omterm{def:b1:euclid:orthogonal}{ortonormal} $(w, n)$ dengan $n = \frac{(a, b)}{\sqrt{a^2+b^2}}$ (yaitu normal
satuan $D$): maka $\norm{M - P_0}^2 = \langle M - P_0,
w\rangle^2 + \langle M - P_0, n\rangle^2$, sehingga $d(M, D) =
\abs{\langle M - P_0, n\rangle}$, dan dengan $aP_{0x} + bP_{0y} =
-c$:
\[
d = \frac{\abs{a x_0 + b y_0 + c}}{\sqrt{a^2 + b^2}} .
\]

\textbf{19.} Dari $E_{\min}/n = v_y - \frac{c_{xy}^2}{v_x}$ dan
$\beta^2 v_x = \frac{c_{xy}^2}{v_x}$:
\[
v_y = \beta^2 v_x + \frac{E_{\min}}n :
\]
yaitu varians total $=$ varians sepanjang garis cocoknya $+$ varians
sisaannya. (Dengan \omterm{def:b1:arith:divides}{membaginya} oleh $v_y$: $1 = \rho^2 + (1 - \rho^2)$, sehingga
bagian varians yang ``dijelaskan'' garisnya adalah $\rho^2$.)

\textbf{20.} Ruang modelnya adalah
$\operatorname{Vect}\bigl(C_1, C_2, C_3\bigr)$ dengan $C_1 = (1)$,
$C_2 = (x_i)$, $C_3 = (x_i^2)$; sehingga meminimalkan $\norm{aC_1 + bC_2 +
cC_3 - b}^2$ menuntun, persis seperti pada pertanyaan 8, ke sistem
Gram $3\times3$, yang matriksnya berentri $\langle C_i, C_j\rangle =
\sum_k x_k^{\,i+j-2}$: yaitu matriks momen yang ditampilkan. Dan ia
dapat dibalik jika dan hanya jika $(C_1, C_2, C_3)$ \omterm{def:b1:vspaces:free}{bebas}
(\cref{exo:b1:euclid:11}); karena sebuah hubungan $a + bx_k + cx_k^2 = 0$
untuk semua $k$ menjadikan setiap $x_k$ akar satu \omterm{def:b1:poly:def}{polinomial}
berderajat $\leq 2$, yang mustahil dengan tiga nilai berbeda kecuali $a
= b = c = 0$. Bandingkanlah matriks momen pada soal akhir pekan
\cref{ch:b1:det}, tempat determinannya berupa kuadrat nilai
Vandermonde.

\textbf{21.} Yang pertama: $r = 2, s = 1, t = 2$, $rt - s^2 = 3 > 0$,
$r > 0$: jadi minimum global yang tegas di titik asalnya. Yang kedua: $r = 2, s
= 3, t = 2$, $rt - s^2 = -5 < 0$: jadi pelana. Yang ketiga: $r = s = t =
2$, $rt - s^2 = 0$: jadi merosot --- tetapi $x^2 + 2xy + y^2 = (x +
y)^2 \geq 0$ lenyap pada seluruh garis $y = -x$: yaitu
minimum global (yang tak tegas) yang tercapai sepanjang sebuah garis, dan tak kasatmata bagi
uji determinannya.

\textbf{22.} Jumlah barunya ($n = 5$): $\bar x = 3.2$, $\bar y =
1.6$, $\sum x_iy_i = 19$ sehingga $c_{xy} = 3.8 - 5.12 = -1.32$,
$\sum x_i^2 = 114$ sehingga $v_x = 22.8 - 10.24 = 12.56$. Kemiringan barunya:
\[
\beta = \frac{-1.32}{12.56} \approx -0.105 :
\]
jadi satu titik mengubah kecenderungan yang jelas naik ($\beta = 1.4$) menjadi
yang sedikit menurun. Adapun galat kuadratnya membebani sebuah sisaan
$\varepsilon^2$, sehingga satu titik yang jauh --- dengan
$\abs{x_{i} - \bar x}$ yang besar \emph{dan} sisaan yang besar --- menguasai
baik $c_{xy}$ maupun $v_x$: jadi \omterm{pb:b1:multivar:1}{kuadrat terkecil} itu mangkus tetapi tak
tegar.

\textbf{23.} Tetapkanlah $W = \sum w_i$ lalu definisikanlah $\bar x_w =
\frac1W\sum w_ix_i$, $\bar y_w$ serupa, $v_x^w = \frac1W\sum
w_ix_i^2 - \bar x_w^2$, $c^w_{xy} = \frac1W\sum w_ix_iy_i - \bar
x_w\bar y_w$. Adapun \omterm{def:b1:logic:map}{pemetaan} $\langle u, v\rangle_w = \sum_i w_iu_iv_i$
merupakan \omterm{def:b1:euclid:def}{hasil kali dalam} pada $\R^n$ (karena $w_i > 0$ memberikan ketegasannya), sehingga
Bagian II berlaku secara harfiah dan \omterm{pb:b1:multivar:1}{persamaan normalnya}, dibagi $W$,
menjadi
\[
\alpha + \beta\bar x_w = \bar y_w,
\qquad
\alpha\bar x_w + \beta(v^w_x + \bar x_w^2) = c^w_{xy} + \bar
x_w\bar y_w ,
\]
dan dari situlah $\beta = c^w_{xy}/v^w_x$ dan $\alpha = \bar y_w -
\beta\bar x_w$: yaitu rumus yang sama, dengan setiap reratanya terbobot.

\textbf{24.} $E_{\min} = \sum\varepsilon_i^2 = 0$ memaksa setiap
$\varepsilon_i = 0$: jadi semua titiknya persis pada garisnya; dan menurut pertanyaan
13 inilah kasus $\abs\rho = 1$. Ujinya: bagi $(1,1), (2,3),
(3,5)$: $\bar x = 2$, $\bar y = 3$, $v_x = \frac{14}3 - 4 =
\frac23$, $c_{xy} = \frac{22}3 - 6 = \frac43$: sehingga $\beta = 2$,
$\alpha = -1$, dan memang $y_i = 2x_i - 1$ bagi ketiga titiknya;
lalu $v_y = \frac{35}3 - 9 = \frac83$ dan $\rho^2 =
\frac{(4/3)^2}{(2/3)(8/3)} = 1$.

\textbf{25.} (i) Bagi fungsi kuadratik penjabaran orde duanya
merupakan sebuah \emph{kesamaan}, sehingga kajian tanda bentuk kuadratiknya
--- yaitu aljabar murni, pertanyaan 1--2 --- menggolongkan titik kritisnya
tanpa teorema Taylor yang diakui. (ii) Adapun \omterm{pb:b1:multivar:1}{persamaan normalnya} berkata
``sisaannya \omterm{def:b1:euclid:orthogonal}{ortogonal} terhadap ruang modelnya'', sehingga minimum
analisisnya \emph{adalah} \omterm{thm:b1:euclid:projection}{proyeksi ortogonal} vektor datanya:
jadi kalkulus dan geometri Euklides menghitung objek yang sama. (iii)
Korelasinya $\rho = c_{xy}/\sqrt{v_xv_y}$ mengukur bagian
$\rho^2$ varians yang dijelaskan garisnya, dan
Cauchy--Schwarz membatasinya: $\abs\rho \leq 1$, dengan kesamaannya hanya
bagi data yang sejajar. (iv) \omterm{def:b1:vspaces:free}{Basis} dan kebebasannya (Bab 18), matriks Gram
dan matriks momen beserta determinannya (Bab 22--23),
dan \omterm{thm:b1:euclid:projection}{proyeksi ortogonal} (Bab 23) semuanya muncul kembali sebagai
bagian kerja satu algoritma. Adapun Bagian II--III mengembangkan
\emph{metode \omterm{pb:b1:multivar:1}{kuadrat terkecil}} (Legendre, Gauss) lewat
\emph{\omterm{pb:b1:multivar:1}{persamaan normalnya}}.

\end{solution}
